\clearpage\pagestyle{otMain}\setlength{\otBlocEspacement}{2.5mm}\fontsize{11}{15}\selectfont\corrigesChapitre{18}{Algorithmique : conditions et variable}
\begin{blocLivre}{cor-18-1}\begin{corrige}{1}{Suivre les valeurs}\demarche{Écrire la valeur de x après chaque instruction.}
x prend d’abord la valeur 4, puis $4+3=7$, puis $7\times2=14$. L’affichage final est donc $\res{14}$. L’ordre des blocs compte.
\label{fig-cor-18-1}\figureGeo[corrige]{script-01}
\end{corrige}\end{blocLivre}
\begin{blocLivre}{cor-18-2}\begin{corrige}{2}{Une entrée utilisateur}\demarche{Traduire la séquence dans son ordre.}
Pour 5 : $(5-2)\times4=\res{12}$. Pour une entrée $x$, le programme donne $\res{4(x-2)=4x-8}$. Le dernier bloc utilise la valeur déjà diminuée de 2.
\label{fig-cor-18-2}\figureGeo[corrige]{script-02}
\end{corrige}\end{blocLivre}
\begin{blocLivre}{cor-18-3}\begin{corrige}{3}{Une boucle d’ajouts}\demarche{Compter les tours et situer l’affichage.}
On part de x = 0 et l’on ajoute 5 à chacun des trois tours : 5, 10, 15. Après la boucle, l’affichage est $\res{15}$. Le bloc « dire » est placé une seule fois, après la répétition.
\label{fig-cor-18-3}\figureGeo[corrige]{script-03}
\end{corrige}\end{blocLivre}
\begin{blocLivre}{cor-18-4}\begin{corrige}{4}{Une condition au seuil}\demarche{Tester aussi la valeur égale au seuil.}
Pour 9, la condition $9>9$ est fausse : le message est « $\res{\text{petit}}$ ». Pour 10, $10>9$ est vraie : le message est « $\res{\text{grand}}$ ». Tester la valeur 9 vérifie la frontière de la condition.
\label{fig-cor-18-4}\figureGeo[corrige]{script-04}
\end{corrige}\end{blocLivre}
\begin{blocLivre}{cor-18-5}\begin{corrige}{5}{Un produit répété}\demarche{Distinguer répétition d’un produit et répétition d’une somme.}
Les valeurs après les tours sont 6, 12, 24 et 48. L’affichage est $\res{48}$. On peut écrire $3\times2^4=48$ ; ajouter quatre fois 2 décrirait un autre programme.
\label{fig-cor-18-5}\figureGeo[corrige]{script-05}
\end{corrige}\end{blocLivre}
\begin{blocLivre}{cor-18-6}\begin{corrige}{6}{Compter des 6}\demarche{Ne modifier la variable que lorsque la condition est vraie.}
La condition est vraie au 2e, au 4e et au 7e tour : compteur vaut $\res3$, le nombre de 6 obtenus. Pour 1 000 lancers, on remplace 10 par 1 000 dans la boucle et l’on fait dire compteur divisé par 1 000 : c’est la $\res{\text{fréquence observée}}$ de 6, proche de $\frac16$ en général.
\label{fig-cor-18-6}\figureGeo[corrige]{script-06}
\end{corrige}\end{blocLivre}
\begin{blocLivre}{cor-18-7}\begin{corrige}{7}{Tracer un carré}\demarche{Additionner les angles de rotation.}
Il trace quatre segments de 30 pixels, avec un quart de tour dans le sens antihoraire entre deux segments. La figure est un $\res{\text{carré de côté 30 pixels}}$. Il revient en $\res{(0\,;0)}$ et retrouve sa direction initiale après un tour complet.
\label{fig-cor-18-7}\figureGeo[corrige]{script-07}
\end{corrige}\end{blocLivre}
\begin{blocLivre}{cor-18-8}\begin{corrige}{8}{Tracer un rectangle}\demarche{Lire tout le contenu répété.}
Chaque tour de boucle trace un côté de 50 pixels puis un côté de 20 pixels, avec des quarts de tour. Deux répétitions ferment un $\res{\text{rectangle de 50 sur 20 pixels}}$. Son périmètre vaut $2(50+20)=\res{140}$ pixels.
\label{fig-cor-18-8}\figureGeo[corrige]{script-08}
\end{corrige}\end{blocLivre}
\begin{blocLivre}{cor-18-9}\begin{corrige}{9}{Deux tests de parité}\demarche{Relier le reste au résultat de la condition.}
Le reste de 18 divisé par 2 est nul : le script dit « $\res{\text{pair}}$ ». Le reste de 19 divisé par 2 vaut 1 : il dit « $\res{\text{impair}}$ ». Les deux branches couvrent tous les entiers.
\label{fig-cor-18-9}\figureGeo[corrige]{script-09}
\end{corrige}\end{blocLivre}
\begin{blocLivre}{cor-18-10}\begin{corrige}{10}{Une touche du clavier}\demarche{Distinguer initialisation et événement répété.}
Chaque appui lance une seule fois le script. Après cinq appuis, score vaut $\res5$. Mettre score à 0 dans ce même script, avant l’ajout, donnerait toujours 1 : l’initialisation doit se faire au début de la partie.
\label{fig-cor-18-10}\figureGeo[corrige]{script-10}
\end{corrige}\end{blocLivre}
\begin{blocLivre}{cor-18-11}\begin{corrige}{11}{Un triangle équilatéral}\demarche{Distinguer l’angle intérieur de l’angle de rotation.}
Les trois côtés mesurent 40 pixels et les trois rotations font $3\times120=360^\circ$ : le lutin revient au départ. La figure est un $\res{\text{triangle équilatéral}}$, d’angles intérieurs $180-120=60^\circ$. Tourner seulement de $60^\circ$ ne refermerait pas la figure en trois côtés.
\label{fig-cor-18-11}\figureGeo[corrige]{script-11}
\end{corrige}\end{blocLivre}
\begin{blocLivre}{cor-18-12}\begin{corrige}{12}{Une interaction simple}\demarche{Repérer le moment où la condition est évaluée.}
Si le contact est présent au moment du test, il dit « $\res{\text{Bonjour !}}$ » ; sinon, il dit « $\res{\text{Je cherche}}$ ». Le script ne contient pas de répétition : il ne fait qu’un seul test à chaque clic sur le drapeau. Une réaction permanente demanderait de répéter le test.
\label{fig-cor-18-12}\figureGeo[corrige]{script-12}
\end{corrige}\end{blocLivre}
\begin{blocLivre}{cor-18-13}\begin{corrige}{13}{Une dimension saisie}\demarche{Relier la variable d’entrée à la figure.}
La boucle trace un carré de côté c. Pour c = 25, la longueur totale vaut $4\times25=\res{100}$ pixels. En général, elle vaut $\res{4c}$ pixels. Il faut saisir une longueur positive.
\label{fig-cor-18-13}\figureGeo[corrige]{script-13}
\end{corrige}\end{blocLivre}
\begin{blocLivre}{cor-18-14}\begin{corrige}{14}{Réparer une boucle}\demarche{Conserver la mémoire d’un tour au suivant.}
À chaque tour, s redevient 0 puis vaut 3 : l’affichage est donc $\res3$. Pour accumuler quatre ajouts, on place « mettre s à 0 » avant la boucle et l’on ne répète que l’ajout. On obtient alors $0+4\times3=\res{12}$.
\label{fig-cor-18-14}\figureGeo[corrige]{script-14}
\end{corrige}\end{blocLivre}
\begin{blocLivre}{cor-18-15}\begin{corrige}{15}{Concevoir et vérifier}\demarche{Tester une valeur de chaque côté du seuil et le seuil lui-même.}
Demander x ; si $3\times x+5>0$, alors dire « positif », sinon dire « nul ou négatif ». Pour 2 : $3\times2+5=11$, $\res{\text{positif}}$. Pour $-2$ : $-1$, $\res{\text{nul ou négatif}}$. Pour $-\frac53$ : 0, $\res{\text{nul ou négatif}}$ ; ce dernier test vérifie le seuil de la condition.
\label{fig-cor-18-15}\figureGeo[corrige]{script-15}
\end{corrige}\end{blocLivre}
